%============================================================
% CH2 — CHAPTER SUMMARY & WHAT COMES NEXT
%   cheat sheet, the five patterns in pictures, recognition table, complexity table,
%   mind map, growth map, self-check quiz (answers at the end of the chapter),
%   readiness checklist, final words
%============================================================

\subsection{The Chapter on One Page}
\label{ssec:ch2-summary-main}

\begin{chaptertheorybox}{Cheat Sheet --- Chapter 2}
\begin{center}
\small
\renewcommand{\arraystretch}{1.45}
\begin{tabularx}{\linewidth}{@{}>{\bfseries}l >{\raggedright\arraybackslash}p{3.4cm} >{\raggedright\arraybackslash}X l@{}}
\toprule
Pattern & What DFS records & Core move & Cost \\
\midrule
P1 Cycles & colours $0/1/2$, \texttt{parent} & edge to a \emph{gray} vertex (undirected: not the parent edge) $\Rightarrow$ cycle $=$ stack from $v$ to $u$ & \VE \\
P2 Subtrees & parent, values per vertex & combine children when \emph{leaving}; pass values down when \emph{entering}; order array for big trees & \bigO{n} \\
P3 Timestamps & $\tin{}$, $\tout{}$, pre-order, sizes & $u$ ancestor of $v \iff \tin{u} \le \tin{v} \wedge \tout{v} \le \tout{u}$; subtree $=$ block of the pre-order & \bigO{n} $+ \bigO{1}$/query \\
P4 Backtracking & the current path & un-mark on exit $\Rightarrow$ every simple path; prune, cap, keep state on the path & exponential \\
P5 Reverse search & reversed edges & one search from the target(s) in $G^R$ instead of one per start & \VE \\
\bottomrule
\end{tabularx}
\end{center}

\textbf{The rhythm:} DFS enters a vertex, disappears into each white neighbour, and leaves; the gray vertices are always the path from the root, and the call stack is that path.

\textbf{The theorems:} parenthesis (intervals nest or are disjoint), white path (descendants $=$ vertices reachable by a white path at entry time), cycles (a directed graph has a cycle iff DFS finds a back edge), undirected DFS has only tree and back edges.

\textbf{The depth rule:} recursion depth $=$ longest DFS path, up to $n$; on a 2\,MB stack it dies near $6.5\cdot10^4$. Iterative DFS with \texttt{it[u]}, or the order array, when $n$ is large.

\textbf{The golden rule:} decide \emph{which hook} the answer lives in --- entering, examining an edge, or leaving --- before writing a line of code.
\end{chaptertheorybox}

\subsection{The Five Patterns in Pictures}

\begin{diagrambox}
\begin{tikzpicture}[
    sv/.style={minimum size=6mm, font=\scriptsize\bfseries},
    pnl/.style={draw=patternborder, fill=patternbg!45, rounded corners=3pt, line width=0.8pt},
    ptitle/.style={font=\scriptsize\bfseries\color{patternborder!80!black}},
    pcap/.style={font=\scriptsize, text width=4.5cm, align=center, text=darkgray}]
    % ---------------- P1 cycles ----------------
    \begin{scope}[shift={(0,4.6)}]
        \draw[pnl] (0,0) rectangle (4.8,4.2);
        \node[ptitle] at (2.4,3.9) {P1 Cycle Detection};
        \node[gvseen, sv] (a) at (0.7,3.0) {1};
        \node[gvsel, sv] (b) at (1.9,3.0) {2};
        \node[gvsel, sv] (c) at (3.1,3.0) {3};
        \node[gvcur, sv] (d) at (4.1,2.1) {4};
        \node[gvdone, sv] (e) at (1.9,1.7) {5};
        \draw[gtree, -{Stealth[length=2mm]}] (a) -- (b); \draw[gtree, -{Stealth[length=2mm]}] (b) -- (c);
        \draw[gtree, -{Stealth[length=2mm]}] (c) -- (d); \draw[gtree, -{Stealth[length=2mm]}] (a) -- (e);
        \draw[gback] (d) to[bend left=25] (b);
        \draw[gcross] (d) to[bend left=10] (e);
        \node[pcap] at (2.4,0.55) {Edge to a gray vertex closes a cycle\\($2,3,4$); to a black one: harmless.};
    \end{scope}
    % ---------------- P2 subtrees ----------------
    \begin{scope}[shift={(5.1,4.6)}]
        \draw[pnl] (0,0) rectangle (4.8,4.2);
        \node[ptitle] at (2.4,3.9) {P2 Subtree Aggregation};
        \node[gvsrc, sv] (r) at (2.4,3.2) {7};
        \node[gvdone, sv] (x) at (1.3,2.2) {3};
        \node[gvdone, sv] (y) at (3.5,2.2) {3};
        \node[gvdone, sv] (x1) at (0.8,1.25) {1}; \node[gvdone, sv] (x2) at (1.8,1.25) {1};
        \node[gvdone, sv] (y1) at (3.0,1.25) {1}; \node[gvdone, sv] (y2) at (4.0,1.25) {1};
        \foreach \a/\b in {r/x, r/y, x/x1, x/x2, y/y1, y/y2} \draw[gtree] (\a) -- (\b);
        \draw[->, thick, primaryblue] (4.5,1.2) -- (4.5,3.2);
        \node[pcap] at (2.4,0.55) {Sizes flow up when a vertex is left:\\$\mathrm{sz}[u] = 1 + \sum \mathrm{sz}[\text{child}]$.};
    \end{scope}
    % ---------------- P3 timestamps ----------------
    \begin{scope}[shift={(10.2,4.6)}]
        \draw[pnl] (0,0) rectangle (4.8,4.2);
        \node[ptitle] at (2.4,3.9) {P3 Entry \& Exit Times};
        \foreach \l/\a/\b/\c in {3/0.4/4.4/primaryblue, 2/0.7/2.3/successgreen, 2/2.6/4.1/successgreen, 1/0.95/1.5/accentorange, 1/2.85/3.4/accentorange}
            \draw[\c, line width=2pt] (\a,\l*0.65+0.6) -- (\b,\l*0.65+0.6);
        \node[font=\tiny\ttfamily] at (2.4,1.15) {( ( ( ) ) ( ( ) ) )};
        \node[pcap] at (2.4,0.55) {Intervals $[\tin{}, \tout{}]$ nest:\\ancestor test in two comparisons.};
    \end{scope}
    % ---------------- P4 backtracking ----------------
    \begin{scope}[shift={(0,0)}]
        \draw[pnl] (0,0) rectangle (4.8,4.2);
        \node[ptitle] at (2.4,3.9) {P4 Backtracking over Paths};
        \node[gvsrc, sv] (s) at (0.6,2.4) {s};
        \node[gv, sv] (p) at (1.8,3.1) {}; \node[gv, sv] (q) at (1.8,1.7) {};
        \node[gv, sv] (w) at (3.0,2.4) {};
        \node[gvdst, sv] (t) at (4.2,2.4) {t};
        \draw[gpath] (s) -- (p); \draw[gpath] (p) -- (w); \draw[gpath] (w) -- (t);
        \draw[ge] (s) -- (q); \draw[ge] (q) -- (w); \draw[ge] (p) -- (q);
        \node[pcap] at (2.4,0.55) {Mark on enter, \textbf{un-mark on exit}:\\every simple path $s \leadsto t$ is tried.};
    \end{scope}
    % ---------------- P5 reverse ----------------
    \begin{scope}[shift={(5.1,0)}]
        \draw[pnl] (0,0) rectangle (4.8,4.2);
        \node[ptitle] at (2.4,3.9) {P5 Reverse Search};
        \node[gvdone, sv] (a1) at (0.7,3.0) {}; \node[gvdone, sv] (a2) at (0.7,1.8) {};
        \node[gvdone, sv] (a3) at (2.0,2.4) {}; \node[gvdead, sv] (a4) at (2.0,1.3) {};
        \node[gvsrc, sv] (tt) at (3.6,2.4) {t};
        \draw[gde] (a1) -- (a3); \draw[gde] (a2) -- (a3); \draw[gde] (a3) -- (tt); \draw[gde] (tt) -- (a4);
        \draw[flow, recursionred] (3.6,3.2) to[bend right=30] node[above, font=\tiny, text=recursionred] {search on $G^R$} (0.9,3.5);
        \node[pcap] at (2.4,0.55) {One search from $t$ along reversed\\edges finds everyone who reaches $t$.};
    \end{scope}
    % ---------------- the three hooks ----------------
    \begin{scope}[shift={(10.2,0)}]
        \draw[pnl, draw=definitionborder, fill=definitionbg] (0,0) rectangle (4.8,4.2);
        \node[ptitle, text=definitionborder] at (2.4,3.9) {Before Any Code: Which Hook?};
        \node[font=\scriptsize, text width=4.3cm, align=left, anchor=north west] at (0.25,3.55) {%
            \textbf{enter $u$:} depth, root-path values, $\tin{u}$, mark on path\\[3pt]
            \textbf{examine $u \to v$:} colour of $v$ $\Rightarrow$ tree / back / forward / cross\\[3pt]
            \textbf{leave $u$:} subtree values, $\tout{u}$, finishing order, un-mark};
        \node[pcap] at (2.4,0.55) {Every pattern fills one or two\\of these three slots.};
    \end{scope}
\end{tikzpicture}
\figcap{The five patterns of the chapter and the one habit that connects them. Orange $=$ current vertex, light orange $=$ gray (on the stack), purple $=$ on the cycle, blue $=$ finished, thick orange $=$ the path being tried.}
\end{diagrambox}

\subsection{Which Pattern? A Recognition Table}

\begin{center}
\small
\renewcommand{\arraystretch}{1.35}
\begin{tabularx}{\linewidth}{@{}>{\raggedright\arraybackslash}p{4.7cm} >{\raggedright\arraybackslash}X c >{\raggedright\arraybackslash}p{2.6cm}@{}}
\toprule
\textbf{The statement says \dots} & \textbf{The graph in disguise} & \textbf{Pattern} & \textbf{Seen in} \\
\midrule
``a round trip'', ``a route back to the start'' & any cycle; undirected: skip the parent, directed: gray target & P1 & Round Trip, Round Trip II \\
``is the order consistent?'', ``contradiction'', ``deadlock'' & constraints as arrows; consistent iff acyclic & P1 & E1.5 \\
``is it a tree / exactly one cycle?'' & counting: $m$ versus $n - c$ & P1 & E1.2, E1.4 \\
``for every vertex, in its subtree \dots'' & post-order aggregation & P2 & White-Black, Highest Score \\
``remove vertex / edge $v$: sizes of the parts'' & $\mathrm{sz}[\text{child}]$ and $n - \mathrm{sz}[v]$ & P2 & Highest Score, E2.2 \\
``add $x$ to every vertex of the subtree of $p$'' & push-down from the root & P2 & Ki \\
``the $k$-th vertex in DFS order of $u$'s subtree'' & pre-order position $+$ size & P3 & Military Problem \\
``descendants of $u$ at depth $d$'' & $\tin{}$ interval $\times$ sorted list per depth & P3 & Count Descendants \\
``remove $k$ cells, keep the rest connected'' & finishing order / a connected prefix & P3 & Maze \\
``all paths'', ``best path'', $n \le 15$ or $\le 25$ cells & backtracking, un-mark on exit & P4 & All Paths, Max Gold \\
``which vertices can reach / are reachable from \dots{} the target'' & reversed graph, one search & P5 & All Ancestors \\
``good ones must escape, bad ones must not'' & block around the bad, search from the exit & P5 & Solve The Maze \\
\bottomrule
\end{tabularx}
\end{center}

\begin{warningbox}[When the table points to the wrong row]
``Shortest'' in the statement means BFS (Chapter~3), not DFS. ``All paths'' with $n = 10^5$ is never backtracking --- it is counting by DP on a DAG (\dpbook{12}{DP on DAGs \& Graphs}). And ``subtree'' questions with \emph{updates} between the queries need the data structures of Chapter~15. Always check the constraints before trusting a row.
\end{warningbox}

\subsection{Complexity at a Glance}

\begin{center}
\small
\renewcommand{\arraystretch}{1.35}
\begin{tabularx}{\linewidth}{@{}l >{\raggedright\arraybackslash}X c c@{}}
\toprule
\textbf{Pattern} & \textbf{Algorithm} & \textbf{Time} & \textbf{Memory} \\
\midrule
P1 Cycles & colour DFS (iterative), parent walk for printing & \VE & \bigO{n+m} \\
P1, yes/no undirected & count components: cycle iff $m > n - c$ & \VE & \bigO{n} \\
P2 Subtrees & one DFS, or BFS order + reverse loop & \bigO{n} & \bigO{n} \\
P3 Timestamps & one DFS; queries by comparisons or binary search & \bigO{n + q \log n} & \bigO{n} \\
P4 Backtracking & all simple paths & up to $\bigO{n!}$ & \bigO{n} \\
P5 Reverse search & one search on $G^R$ (per target, or multi-source) & \VE & \bigO{n+m} \\
\bottomrule
\end{tabularx}
\end{center}

\noindent Rule of thumb for P4: count the paths before coding. A $6 \times 6$ grid graph already has $1\,262\,816$ simple corner-to-corner paths (Lesson~\ref{ssec:ch2-toolbox}); $15$ vertices of a DAG have at most $2^{13} = 8192$ paths from the first to the last.

\subsection{Mind Map}

\begin{diagrambox}
\begin{tikzpicture}[
    mindmap, grow cyclic, every node/.style={concept, execute at begin node=\hskip0pt},
    concept color=primaryblue!70, text=white, font=\scriptsize\bfseries,
    level 1/.append style={level distance=4.1cm, sibling angle=72, font=\scriptsize\bfseries},
    level 2/.append style={level distance=2.3cm, sibling angle=40, font=\tiny}]
\node[font=\small\bfseries, minimum size=2.6cm] {DFS}
    child[concept color=successgreen!80!black] { node {Mechanics}
        child { node {call stack} }
        child { node {iterative it[u]} }
        child { node {depth limits} } }
    child[concept color=accentorange] { node {Timestamps}
        child { node {colours} }
        child { node {tin / tout} }
        child { node {parenthesis} } }
    child[concept color=softpurple] { node {DFS tree}
        child { node {tree, back} }
        child { node {forward, cross} }
        child { node {white path} } }
    child[concept color=patternborder] { node {5 Patterns}
        child { node {cycles, subtrees} }
        child { node {times, paths} }
        child { node {reverse} } }
    child[concept color=warningred!80!black] { node {Traps}
        child { node {gray vs visited} }
        child { node {parent edge} }
        child { node {stack overflow} } };
\end{tikzpicture}
\end{diagrambox}

\subsection{How Each Pattern Grows Through the Book}

\begin{diagrambox}
\begin{tikzpicture}[x=3.15cm, y=1.25cm,
    pat/.style={rectangle, rounded corners=3pt, draw=patternborder, fill=patternbg, line width=1pt,
                font=\scriptsize\bfseries, align=center, text width=2.2cm, minimum height=9mm},
    lvl/.style={rectangle, rounded corners=2pt, draw=#1, fill=#1!10, font=\tiny, align=center,
                text width=2.45cm, minimum height=8mm}]
    \node[font=\scriptsize\bfseries, color=deepblue] at (0,0.8) {Chapter 2};
    \node[font=\scriptsize\bfseries, color=successgreen] at (1,0.8) {Part I (Ch.~3--6)};
    \node[font=\scriptsize\bfseries, color=softpurple] at (2,0.8) {Part II--III (Ch.~7--17)};
    \node[font=\scriptsize\bfseries, color=warningred] at (3,0.8) {Part IV--VI (Ch.~18--35)};
    \node[pat] (p1) at (0,0)  {P1 Cycles};
    \node[pat] (p2) at (0,-1) {P2 Subtrees};
    \node[pat] (p3) at (0,-2) {P3 Timestamps};
    \node[pat] (p4) at (0,-3) {P4 Backtracking};
    \node[pat] (p5) at (0,-4) {P5 Reverse};
    \node[lvl=successgreen] (a1) at (1,0)  {topological order, cycle reporting (Ch.~5)};
    \node[lvl=successgreen] (a2) at (1,-1) {components as units (Ch.~4)};
    \node[lvl=successgreen] (a3) at (1,-2) {finishing order = topological order (Ch.~5)};
    \node[lvl=successgreen] (a4) at (1,-3) {BFS over states instead (Ch.~3)};
    \node[lvl=successgreen] (a5) at (1,-4) {multi-source BFS (Ch.~3)};
    \node[lvl=softpurple] (b1) at (2,0)  {negative cycles (Ch.~8)};
    \node[lvl=softpurple] (b2) at (2,-1) {tree fundamentals, rerooting (Ch.~13)};
    \node[lvl=softpurple] (b3) at (2,-2) {Euler tour + range queries (Ch.~13, 15)};
    \node[lvl=softpurple] (b4) at (2,-3) {DP on DAGs, bitmask DP (DP book)};
    \node[lvl=softpurple] (b5) at (2,-4) {reversed Dijkstra (Ch.~7)};
    \node[lvl=warningred] (c1) at (3,0)  {low-links: bridges, cut vertices (Ch.~18)};
    \node[lvl=warningred] (c2) at (3,-1) {small-to-large, centroids (Ch.~15, 17)};
    \node[lvl=warningred] (c3) at (3,-2) {SCC by finishing times (Ch.~19)};
    \node[lvl=warningred] (c4) at (3,-3) {Hamiltonian paths on small graphs (Ch.~23)};
    \node[lvl=warningred] (c5) at (3,-4) {Kosaraju's second pass on $G^R$ (Ch.~19)};
    \foreach \i in {1,...,5}{\draw[flow] (p\i) -- (a\i); \draw[flow] (a\i) -- (b\i); \draw[flow] (b\i) -- (c\i);}
\end{tikzpicture}
\figcap{Every arrow means ``the same DFS with one more piece of bookkeeping''. Bridges and strongly connected components, the two most famous DFS algorithms, are built entirely from this chapter's timestamps and back edges.}
\end{diagrambox}

\subsection{Self-Check Quiz}

Answer without looking back.
\begin{enumerate}[leftmargin=*]
    \item What is on the call stack of a recursive DFS at any moment, in graph terms?
    \item Why does the ``push all neighbours'' iterative DFS give no exit times?
    \item State the parenthesis theorem. What does $\tin{u} < \tin{v} < \tout{v} < \tout{u}$ tell you?
    \item Name the four edge types of a directed DFS and how each one is recognised when the edge is examined.
    \item Why does an undirected DFS have neither forward nor cross edges?
    \item In a directed graph, why is ``the target is already visited'' not a cycle test? Give a 3-vertex example.
    \item In an undirected multigraph, why must the parent \emph{edge} be skipped rather than the parent vertex?
    \item An undirected graph has $n$ vertices, $m$ edges and $c$ components. When does it contain a cycle?
    \item Give the formula for the sizes of the parts left when a vertex $v$ is deleted from a rooted tree.
    \item With pre-order numbers $\mathrm{pos}$ and sizes $\mathrm{sz}$, which positions does the subtree of $u$ occupy?
    \item What exactly changes between DFS and backtracking, and what does it cost?
    \item How do you find all vertices that can reach $t$ with a single search?
\end{enumerate}

\noindent\textit{Answers are at the end of the chapter, after the solution sketches (page~\pageref{ssec:ch2-quiz-answers}).}

\subsection{Ready for Chapter 3?}

\begin{checklistbox}[Move on when every box is ticked]
\begin{itemize}[leftmargin=*]
    \item[$\square$] I can write recursive DFS, iterative DFS with \texttt{it[u]} and the order-array loop from memory, and I know when recursion is too deep.
    \item[$\square$] I can compute $\tin{}$, $\tout{}$, pre-order and subtree sizes, and use the ancestor test.
    \item[$\square$] I can classify any edge of a DFS as tree, back, forward or cross, by hand and in code.
    \item[$\square$] I can find and print a cycle in directed and undirected graphs, including multigraphs and self-loops.
    \item[$\square$] I can solve a ``for every subtree'' problem on $2\cdot10^5$ vertices without recursion.
    \item[$\square$] I can recognise when backtracking is affordable, and when a search should be run backwards.
    \item[$\square$] I have solved at least 20 of the practice problems of Tiers 1--3, including at least 4 from Tier 3.
\end{itemize}
\end{checklistbox}

\subsection{What Comes Next}

DFS answers ``what is the structure?'': cycles, subtrees, ancestors, all the paths. It does \emph{not} answer ``how far?'' --- the path DFS finds to a vertex can be absurdly long (Lesson~\ref{ssec:ch2-toolbox}). Chapter~3 (BFS) explores in layers instead, and every vertex is reached by a shortest path; most of the ``minimum number of moves'' problems of contests live there. Chapters~4--6 then combine both searches into components, bipartite graphs, topological orders and disjoint-set unions. Later, Chapter~18 (bridges and articulation points) and Chapter~19 (strongly connected components) return to the DFS tree of this chapter and squeeze out much more of its structure with one extra number per vertex.
